\chapter{Discussion}

\section{Assessment of the hypotheses}

\textbf{H1 --- supported as to structure and volume.} The public data make it
possible to reconstruct sequences with the same structure as Lee's, in sufficient
volume. Their statistical content, however, differs from that of equities: persistent
direction, multi-tick jumps and times discretized by the grid
(Section~\ref{sec:caracteristicas}).

\textbf{H2 --- partially supported.} The LSTM's error is of the same order as the
MLE's: at the fixed point, 0.0111 versus 0.0103; on the synthetic test, the median
per-path loss is practically equal. The substantially lower latency was not confirmed
at batch size one against an efficient MLE (Section~\ref{sec:custo}); it holds for
batched inference.

\textbf{H3 --- not supported on synthetic data.} The Transformer did not reduce the
error: its mean macro MSE was higher than the LSTM's, with no significance across
seeds, and the Transformer ensemble lost to the LSTM ensemble in every regime. Nor
did it perform better in the near-critical regime, where long memory would be most
relevant. Its advantage was latency on MPS (5 ms versus 30 ms for the LSTM), at the
cost of 2.7 times longer training. The result holds for the configuration tested;
there was no hyperparameter search, and the \emph{post-LN} Transformer was trained
without learning-rate warm-up. On real windows, the fine-tuned architectures are indistinguishable on the test block, and in the October 2023
window the Transformer has a lower loss than the selected LSTM in 64\% of the windows
(Section~\ref{sec:real-robustness}). This robustness under a change of regime is the first
evidence in favor of H3, but it comes from a post hoc comparison and not from the
pre-registered selection, so it is a hypothesis for future work rather than a confirmation.
The addendum (Section~\ref{sec:addendum}) sharpens it. On a new week, pre-registered and read
once, no Transformer arm beat the frozen LSTM, while the average of the two networks had the
lowest loss; in the $\beta$ sweep, the LSTM errs less inside the training range and the
Transformer above it. This complementarity depends on the prior: pretrained on a prior built from the real labels
(Section~\ref{sec:addendum-scope-b}), the LSTM alone is the best estimator on March and
October, and averaging no longer helps. The Transformer keeps one advantage, for volatility:
it preserves more of the ranking of the windows (Section~\ref{sec:addendum-volatility}).

\section{Why the networks win on aggregate MSE}

The MLE's macro MSE (0.506) is about 25 times that of the networks, but this
difference stems from two mechanisms, not from greater accuracy of the networks on a
typical path. The first is identifiability (Section~\ref{sec:identificabilidade}): in
the low-\(g\) regime, the MLE produces extreme values of \(\beta\) that legitimately
maximize the likelihood, and 1\% of the paths account for 94\% of the squared error
of \(\beta\). The second is the prior (Section~\ref{sec:bayes}): the networks
approximate the posterior mean under the training prior and never leave the interval
\([0.8,\,2.0]\) for \(\beta\), which bounds their worst-case error. This robustness is
useful, but it comes at the cost of shrinkage bias and depends on the prior
representing the data.

\section{Comparison with Lee (2023)}

Lee's qualitative conclusions hold: the network reaches error of the same order as
the MLE, with shrinkage bias, and the MLE is nearly unbiased when the parameter is
well identified. The absolute MSE values are not comparable across the two works,
since they depend on the parameter ranges, which the paper does not publish. The
results diverge from the paper in two respects: the speed gain of two orders of magnitude is not
reproduced with an efficient MLE (a likelihood computed by double sum over pairs of events
recovers it; Section~\ref{sec:custo}), and evaluating by aggregate MSE alone hides the
identifiability failures that initialization at the true value tends to mask.

\section{What the real data show}

\paragraph{The model.} In 99.8\% of the calibration windows the Hawkes model is significantly better than a
homogeneous Poisson process, so the clustering of price changes is real and large
(Section~\ref{sec:real-diagnostics}). The model is, however, misspecified: the rescaled
residuals keep serial dependence in about half of the windows, against 5\% in simulated
controls observed on the same grid. A single exponential kernel describes most of the
clustering, and the 100~ms grid explains only part of the
departure from \(\mathrm{Exp}(1)\). On real data, therefore, the networks learn
the MLE of an approximating model. It is also the MLE of the process seen through
the grid: merging events in the same cell slows the apparent decay, and the grid label is
about 58\% of the continuous-time $\beta$ (Section~\ref{sec:addendum-grid}). A network trained on
the true parameters of gridded paths undoes this bias in simulation, better on average than
inverting the map, but on real windows the two corrections disagree
(Section~\ref{sec:addendum-grid-correction}).

\paragraph{The parameters.} The estimated parameters (\(\alpha_{\mathrm{cross}}\approx0\),
\(\beta\approx2.7\)) lie outside the support of the synthetic prior. The null
cross-excitation contrasts with the typical point for equities, where it dominates (at Lee's
fixed point, \(\alpha_{\mathrm{cross}}>\alpha_{\mathrm{self}}\)), and is consistent with the
directional persistence measured in the data. Applied as they are, the synthetic networks
shrink their estimates into the prior and are worse than the constant mean of the labels.

\paragraph{The networks.} Trained on MLE labels of real windows, as \citet{lee2023} suggests,
the networks become useful only when they start from the synthetic weights: from scratch,
1,063 windows are too few: on the test block the LSTM loses to the label mean and the
Transformer beats it only modestly. Fine-tuned, both explain most of the label variance,
with test macro MSE five times lower than the label mean. Synthetic training thus works as
an informative initialization even when its prior does not cover the data. The LSTM chosen on
validation and the fine-tuned Transformer cannot be told apart on the test block; in October
the Transformer degrades less. The selected LSTM was also undertrained: with a 1,000-epoch cap and
nothing else changed, its loss on a new week falls by a fifth (Section~\ref{sec:addendum-holdout}).
The prior had a larger effect than the budget: pretrained on a prior built from the real labels and
simulated on the same grid (a reduced scope B), the LSTM cuts the loss on a second new week by
almost two thirds, mostly because it extrapolates $\mu$ to activity levels that the January
calibration did not cover (Section~\ref{sec:addendum-scope-b}).

\paragraph{The volatility.} With the MLE, the Hawkes volatility ranks the windows like the
realized volatility (Spearman 0.78 on the test block at \(\Delta=10\)~s) and is about 1.27
times higher at horizons of 5 to 30~s. The ratio falls from about 1.6 at 0.1~s to 1.27 at
10~s, so the realized variance per unit time grows with \(\Delta\), as expected when
successive price changes are positively autocorrelated, which is the directional persistence
that the self-excitation describes. The Hawkes expression is a long-run quantity
(\(\Delta\to\infty\)); since the ratio is still about 1.3 at 30~s, either the realized
variance has not yet reached its long-run level or the model overstates the persistence. With the network,
the level is kept but the ranking falls (0.28): the network shrinks its estimates toward the
mean of the training labels, which compresses the differences between windows. For
volatility, the estimator must also preserve the cross-section of the windows. The Transformer does this better than the LSTM
(Spearman 0.69 against 0.47 on February, Section~\ref{sec:addendum-volatility}), but the MLE
remains the best input (0.77).

\section{Threats to validity}

\begin{itemize}
  \item \textbf{Internal:} leakage through overlapping windows; tuning after observing
        the test set. Mitigated by chronological splits and a test set locked behind a
        flag.
  \item \textbf{Construct:} on real data the label is the MLE of a model that the
        residual diagnostics reject in part, so the networks are evaluated on reproducing
        an approximating fit; $\alpha_{\mathrm{cross}}$ lies on the boundary of the
        parameter space, where the usual asymptotic theory does not hold. The labels are
        also biased by the 100~ms grid (about 58\% of the continuous-time $\beta$), which
        jittering inside the cell does not undo.
  \item \textbf{Prior:} the synthetic metrics depend on the sampled ranges, which were
        not calibrated on the data; the reduced scope B, which calibrates them, was run for
        the LSTM only, with a single prior design.
  \item \textbf{External:} two assets and two windows do not represent the whole
        market; the empirical test block covers three days, including a weekend.
  \item \textbf{Statistical:} five seeds quantify only part of the training
        variability; the LSTM showed two convergence levels across seeds, and one
        seed reached the 300-epoch limit while still improving. On real windows, the
        fine-tuned LSTM reached its best epoch at 291--299 in every seed, so it was also
        limited by the epoch budget; the addendum confirms that a longer budget lowers its
        loss on a new week. The comparison between fine-tuned networks on the test
        and October blocks is post hoc.
  \item \textbf{Reproduction:} the paper omits ranges, seeds and training details.
  \item \textbf{Computational:} latency depends on hardware, batch size, warm-up and
        contention between processes.
\end{itemize}
